\chapter{Evaluation und Vergleich}\label{chap:Evaluation-und-Vergleich}

Dieses Kapitel führt die Klassifikationsergebnisse und die XAI-Auswertung zusammen. Entscheidend ist nicht nur, welches Modell die höchsten Kennzahlen erreicht. Ebenso wichtig ist, ob die Fehlerstruktur und die verwendeten Merkmale nachvollziehbar bleiben.

\section{Bewertung der Modellleistung}

Die vier untersuchten Modelle erreichen im Testsetup sehr ähnliche Ergebnisse. Das MLP erzielt mit einer Accuracy von 0{,}923 und einem Macro-F1 von 0{,}934 die höchsten Werte. HistGradientBoosting erreicht 0{,}920 und 0{,}932. Random Forest erreicht 0{,}918 und 0{,}931, die logistische Regression 0{,}917 und 0{,}931. Zwischen dem besten und dem schwächsten Modell liegen nur 0{,}006 Accuracy-Punkte und 0{,}004 Macro-F1-Punkte.

Aus den reinen Kennzahlen lässt sich deshalb kein klar überlegenes Modell ableiten. Das MLP ist im Testlauf knapp vorn. Die logistische Regression bleibt als lineares Vergleichsmodell wichtig, weil sie trotz geringerer Komplexität eine ähnliche Leistung erzielt.

\section{Fehlerstruktur der Modelle}

Die Konfusionsmatrizen zeigen, dass die Fehler nicht zufällig über alle Klassen verteilt sind. \texttt{BOMBAY} wird sehr zuverlässig erkannt. Das passt zur fachlichen Erwartung, weil diese Sorte im Datensatz deutlich größere Bohnen aufweist und sich dadurch stärker von den übrigen Klassen abhebt.

Schwieriger sind die Trennungen zwischen \texttt{DERMASON} und \texttt{SIRA} sowie zwischen \texttt{BARBUNYA} und \texttt{CALI}. Diese Fehlermuster treten bei mehreren Modellen auf. Das spricht dafür, dass die betroffenen Sorten mit dem reduzierten Feature-Set schwerer zu unterscheiden sind.

\section{Einordnung der XAI-Ergebnisse}

Die XAI-Auswertung umfasst alle vier finalen Modelle. PFI, globale SHAP-Werte sowie lokale SHAP- und LIME-Erklärungen wurden auf Logistic Regression, Random Forest, HistGradientBoosting und MLP angewandt. Über die Modelle hinweg treten \texttt{Area} und \texttt{ShapeFactor1} wiederholt unter den höchsten Rängen auf. Je nach Modell ergänzen \texttt{ShapeFactor2}, \texttt{Roundness}, \texttt{Compactness} oder \texttt{AspectRatio} dieses Muster. Das passt dazu, dass Bohnensorten über sichtbare Unterschiede in Größe, Rundheit und Form beschrieben werden können.

PFI bewertet den Leistungsverlust nach dem Vertauschen eines Merkmals. SHAP betrachtet Beiträge zur Vorhersage. LIME erklärt eine lokale Umgebung um eine einzelne Beobachtung. Deshalb unterscheiden sich einzelne Rangplätze, obwohl das Gesamtbild ähnlich bleibt. Mehrere Features sind stark korreliert. Bei solchen Merkmalen kann sich Wichtigkeit zwischen ähnlichen Variablen aufteilen.

\section{Vergleich von Leistung und Erklärbarkeit}

Das beste Modell nach Macro-F1 ist das MLP. Die XAI-Auswertung beschränkt sich aber nicht auf das beste oder auf ein besonders gut erklärbares Modell. Sie vergleicht alle vier Modellfamilien und vermeidet damit die Schlussfolgerung, dass die Befunde des Random Forest automatisch für MLP, HistGradientBoosting oder logistische Regression gelten.

Der Vergleich zeigt einen praktischen Zielkonflikt. Komplexere Modelle können leicht bessere Kennzahlen erreichen, liefern aber nicht automatisch leichter berechenbare Erklärungen. Für das MLP war deshalb ein begrenzter modellagnostischer SHAP-Ansatz nötig. Die Ergebnisse bleiben für Rangfolgen innerhalb des MLP aussagekräftig. Absolute SHAP-Werte werden wegen der verschiedenen Explainer nicht modellübergreifend gleichgesetzt.
